Un skieur aborde une piste constituée de trois parties (voir figure).

\begin{figure}[H]
    \centering
    \begin{tikzpicture}[thick]
        \def\be{30}
        \def\th{30}
        \def\AB{4}
        \def\BC{8}
        \def\r{2.4}
        \draw[very thick] (0,0) node[point,label=$A$] (A) {} --
            ++(-\be:\AB) node[point,label=$B$] (B) {} --
            ++(\BC,0) node[point,label=$C$] (C) {} arc(90:0:\r)
            node[point,label=right:$D$] (D) {};
        \draw[densely dashed,thick]
            (B.center) -- (B.center -| A.center);
        \draw ([xshift=-1cm]B.center) arc (180:180-\be:1)
            node[pos=0.6,left] {$\beta$};
        \draw[thick,densely dashed] (C) -- ++(0,-\r) coordinate (O)
            node[below left,inner sep=1pt] {$O$} -- (D);
        \draw[thick,densely dashed] (O) -- ++(\th:\r)
            node[point,label=above right:$M$] {};
        \draw ([xshift=1cm]O.center) arc (0:\th:1)
            node[pos=0.6,right] {$\theta$};
    \end{tikzpicture}
\end{figure}

Le skieur, de masse $m=\qty{80}{\kg}$, part du point $A$ à une vitesse
$v_{A}=\qty{3}{\m\per\s}$.

\textbf{On donne~:} $g=\qty{10}{\N\per\kg}$.
\begin{enumerate}
    \item La première partie $AB$, de longueur $AB=\qty{4}{\m}$, est un plan
          incliné d'angle $\beta=\ang{30}$ sur l'horizontal.
    
          Les frottements sont négligeables sur la partie $AB$.
          
    
          \begin{enumerate}
              \item Faire le bilan et représenter les forces qui s'exercent sur
                    le skieur sur la partie $AB$.~\textbf{(1pt)}
              \item En appliquant le théorème de l'énergie cinétique entre $A$
                    et $B$, montré que la vitesse du skieur en $B$ est
                    $v_{B}=\qty{7}{\m\per\s}$.~\textbf{(1,5pt)}
          \end{enumerate}
    
    \item La partie $BC$, horizontale, de longueur $BC=\qty{8}{\m}$~; les
          frottements sont équivalents à une force d'intensité
          $f=\qty{120}{\N}$.
          \begin{enumerate}
              \item Faire le bilan et représenter les forces qui s'exercent sur
                    le skieur sur la partie $BC$.~\textbf{(1,5pt)}
              \item Calculer sur la partie $BC$, le travail de la force de
                    frottement $\vec{f}$.~\textbf{(1pt)}
              \item En appliquant le théorème de l'énergie cinétique entre $B$
                    et $C$, calculer la vitesse $v_{C}$ du skieur en
                    $C$.~\textbf{(1,5pt)}
          \end{enumerate}
    \item La partie $CD$, est un arc de cercle de centre $O$ et de rayon
          $r=\qty{2.4}{\m}$. Les frottements sont négligeables sur la partie
          $CD$. La position de point $M$ est repéré par l'angle
          $\theta=(\widehat{\overrightarrow{OD},\overrightarrow{OM}})$
          \begin{enumerate}
              \item Exprimer la vitesse $v_{M}$ du skieur au point $M$ en
                    fonction de l'angle $\theta$, le rayon $r$ et
                    $V_{C}$.~\textbf{(1pt)}
              \item Calculer l'angle $\theta$ pour que la vitesse du skieur au
                    point $M$ soit $v_{M}=\qty{7}{\m\per\s}$.~\textbf{(1pt)}
          \end{enumerate}
\end{enumerate}

